%/*% pdflatex -file-line-error -interaction nonstopmode %

\documentclass{article}
\usepackage{amsmath}
\usepackage{amsfonts}
\usepackage{hyperref}
\usepackage{exercises}
\usepackage{amsthm}
\usepackage{yhmath}
\newcommand*{\oldneg}{\mathord{\sim}}

\theoremstyle{definition}
\newtheorem{definition}{Definition}

\setlength{\parindent}{0pt}
\setlength{\parskip}{10pt}

\title{Chapter 13 Self Notes}
\begin{document}
\maketitle

\section*{Necessity of the exclusion of free occurences in $C$}
One may ask, why does $C$ necessarily need to have no free occurences of $x$ (in the $Group\ Inference\ Rules)$ to make sense in relation to $FOL$?
 Well, if we did not, we very well could have a scenario in which $C\supset F(x)$ is provable, and $C\supset F(t)$ is provable for some $t$ with both 
$C$, $F(t)$ being provable, and also some $t$ such that both $C$ and $F(t)$ are refutable. Hence, $\forall xF(x)$ would not be true (in $FOL$), and 
in the case where $C$ is true, we would have $C\supset \forall xF(x)$ being not true. We can use similar reasoning for the second inference rule.\newline
\newline A similar question arises for the need of substituting a term $t$ for for $x$ in $F$ in the $Group\ II\ Axioms$ such that $x$ is not bound by any 
variables in $t$. Well, it could be the case that 
$t$ is some variable $y$ which is bounded in $F$, and such that $y$ is not free for $x$ in $F$. In this case, substituting $y$ for $x$ would essentially 
'add' a new relation that $y$ is quantified over, as now we would have $y$ being placed in the condition $x$ was previously 
placed in. Well, this scenario does not obey the laws of $FOL$, and thus we should avoid it.

\section*{Meaning of 'The same as'}
We should note that, when Smullyan says being dyadically enumerable is the same as being $\Sigma_1$, it does not mean these conditions are equivalent. 
Instead, he means that dyadic enumerability implies being $\Sigma_1$. And so, when he says that we have shown Peano Arithmetic to be a formal system, 
he really means that we have shown $P_0$ and $R_0$ to be $\Sigma_1$, and this is all we need to apply the Gödel and Rosser theorems for formal 
systems (as, in those proofs, we only use the fact that $P_0$ and $R_0$ being dyadically enumerable means they are also $\Sigma_1$).

\section*{Necessity of Using Sentences Instead of Formulae}
When going through these last four chapters, one will inevitably ask the question, why are all of the sets being used defined for sentences, not for 
formulas? To start, let us consider the system of First-order arithmetic. Well, this system deals solely with the notion of truth. But truth is in fact 
not defined for sentences, and therefores, if we were to change the sets to be defined for formulas, it would actually not change anything from them 
being defined only for sentences (and this holds true for whenever sets are defined that deal with truth). Now, in Chapter 10, we are talking about 
systems solely dealing with provability. We are given a well defined set of provable $sentences$, but no guarantee that we have a well defined set of 
provable $formulas$. Of course, there is a well defined set of provable $formulas$ in Peano Arithmetic, but the reason we define all the sets of the 
results of Gödel and Rosser only with sentences is so that we only $need$ to know that there is a well defined set of sentences (which we have in Peano 
Arithmetic), which means that getting all these results with sets defined with sentences is a more general result than what could be obtained with 
formulas (in the sense that these are formulas and sentences of FOL). 
\end{document}

